\documentclass[conference]{IEEEtran}
\IEEEoverridecommandlockouts

\usepackage{cite}
\usepackage{amsmath,amssymb,amsfonts}
\usepackage{graphicx}
\usepackage{booktabs}
\usepackage{url}
\usepackage{xcolor}
\usepackage{multirow}
\usepackage{enumitem}
\usepackage{microtype}
\usepackage[hidelinks]{hyperref}

\newcommand{\codebleu}{\textsc{CodeBLEU}\xspace}
\usepackage{xspace}
\newcommand{\passatk}{\texttt{pass@k}\xspace}
\newcommand{\passat}[1]{\texttt{pass@#1}\xspace}
\newcommand{\compileatk}{\texttt{compile@k}\xspace}
\newcommand{\compileat}[1]{\texttt{compile@#1}\xspace}
\newcommand{\pp}{\,pp\xspace}
\newcommand{\fnz}{\texttt{fn0}\xspace}

\begin{document}

\title{The Identifiers Are the Specification:\\
Auditing What Neural Decompilation Benchmarks Measure}

\author{\IEEEauthorblockN{Raafat Abualazm, Ayman Abo Elhassan, and Amr G. Wassal}
\IEEEauthorblockA{\textit{Computer Engineering Department, Faculty of Engineering} \\
\textit{Cairo University}, Giza, Egypt \\
raafat.202210476@eng-st.cu.edu.eg}}

\maketitle

\begin{abstract}
Neural decompilers are evaluated on benchmarks ported from HumanEval. An audit
of six published evaluation configurations finds that where the original
identifiers reach the model is decided by the benchmark split rather than the
system, and that none of the six reports an identifier-scrubbed control arm. We
ask what such scores measure by decomposing the
model's input on a frozen checkpoint, varying one channel at a time and
replicating the channel arms across five sampling seeds, extended to six under a
pre-registered amendment. Arms outside that family carry fewer seeds and are
labelled as such.

The channels separate. On the leakage-clean 174-task set across six sampling
seeds, a typed contract supplied at inference with \emph{no training} moves
\passat{10} from $4.8$ to $10.2$ (mean paired difference $+5.3$) and
\compileat{10} from $126.0$ to $168.2$ ($+42.2$); both effects are positive in
every one of the six seed blocks, so an exact paired sign-flip test gives
$p{=}2/2^6=0.031$. The arms share seeds, so this blocked test is the valid one,
and six blocks is the fewest it can resolve. The other two channels are
directional at this precision: deleting recovered constants moves \passat{10} by
$-2.2$ ($p{=}0.25$) with no compilation effect detected, and permuting which
binary belongs to which task moves \compileat{10} by $-8.7$ (sign-consistent,
$p{=}0.031$) with no correctness effect detected ($+0.4$, $p{=}0.875$). A
sixth seed was added under a pre-registered amendment with a single confirmatory
test and a hard stop; it is the only inference we claim at $\alpha{=}0.05$. A pre-registered oracle control that supplies
the same typed contract but replaces the binary payload with the empty string ---
an out-of-distribution input --- reaches the highest compilation we measured,
$173.5$ of $174$, while solving no more tasks than baseline.

The dominant channel is the bundle of semantic identifiers. On a frozen
checkpoint, replacing the function name, every parameter name and the symbol in
the disassembly with opaque equivalents --- leaving types, arity and the
instruction stream byte-identical --- \emph{improves} compilation by $7.5$\pp
while collapsing \passat{10} from $27.27\%$ to $4.55\%$, losing $35$ tasks and
gaining none (exact $p{\approx}6{\times}10^{-11}$). Those identifiers function as
a compact specification; our design does not separate recall of a memorised
solution from inference over a descriptive name, and we claim only the former
where the evidence supports it. Independently, an untrained base model prompted
with the signature alone, never shown a binary, is not measurably worse than the
same model fine-tuned on binaries ($-0.0195$, $p{=}0.549$).

We also establish what this protocol can resolve at all. Re-running one frozen
checkpoint under five seeds moves \passat{10} across $2$--$6$ of $174$ tasks
($\mathrm{SD}=1.79$), so differences under about three tasks --- including
several in this literature and several of our own --- are indistinguishable from
reseeding the same model. A stronger form of the same defect is that \passatk can
be reproducible in aggregate while being irreproducible in composition: our
strongest arm reports $10.0\pm1.0$ solved tasks across five seeds, yet those five
runs solve $24$ distinct tasks and agree on only $4$ (mean pairwise Jaccard
$0.375$). The same runs scored on compilation, which operates an order of
magnitude higher, reproduce at $0.973$. Against that floor, an exploratory sweep
of sixteen architectural interventions and a rejection-sampling programme across
six checkpoints yield nothing that survives correction, and a frontier reasoning
model with ten attempts solves four tasks.

We argue that a name-scrubbed arm, a seed-replicated noise floor and an
input-channel decomposition are together necessary for a decompilation result to
be interpretable, and we document an evaluation-leakage defect that invalidated
part of our own earlier evidence.
\end{abstract}

\begin{IEEEkeywords}
neural decompilation, benchmark validity, negative results, evaluation
methodology, reverse engineering, large language models
\end{IEEEkeywords}

\section{Introduction}
Decompilation recovers readable source from compiled binaries. Traditional
decompilers emit correct but unidiomatic output with synthesised identifiers, and
a growing literature applies large language models to close that
gap~\cite{llm4decompile}, reporting functional-correctness rates that look
encouraging.

Nearly all of that literature evaluates on benchmarks ported from
HumanEval~\cite{chen2021evaluating}. This paper asks what those scores measure.
Our answer, from a controlled decomposition of the model's input, is that they
substantially measure whether the model can identify the task from the human
identifiers it is given --- which together form a compact natural-language
specification --- and that the remaining channels divide cleanly between those
that produce compilable code and those that produce correct code.

We are deliberate about the mechanism we do \emph{not} claim. A descriptive
identifier such as \texttt{sortViolationsBySeverity} may be recalled from
pretraining or read as a specification and solved afresh; our manipulation
removes the identifier and cannot distinguish these. We therefore describe the
defect as \emph{specification leakage} and reserve claims of memorisation for
where separate evidence supports them.

\subsection{Which evaluations this critique applies to}
\label{sec:scope}
We surveyed prompt-construction code, released datasets and prompt templates
rather than assert where the practice occurs. The survey covers published neural
decompilers and decompilation benchmarks from 2024--2026 for which a prompt
template, prompt-construction code or a released evaluation dataset is publicly
available; a configuration whose prompt could not be established from an artifact
was recorded as unverifiable and excluded rather than inferred.
Table~\ref{tab:survey} reports the six configurations over four systems that met
that bar, each row checked against the artifact independently of the search that
located it.

The finding is that \textbf{identifier exposure is a property of the benchmark
split and the input channel, not of the system}. No verified C-target system
supplies a name, parameter names or a signature as a dedicated prompt
\emph{field}; LLM4Decompile, SALT4Decompile and SK2Decompile share one
single-slot template. But the symbol reaches the model inside the assembly
regardless. On HumanEval-Decompile that symbol is the semantically empty
\texttt{func0} in $656/656$ rows, so nothing leaks, and the benchmark should be
credited for it. On the \emph{same system's} GitHub2025 and ExeBench splits the
first line of the assembly is the true symbol ---
\texttt{mi\_cmp\_dynamic\_unique:}, \texttt{StateIdle:} --- and C++ rows carry the
demangled signature with parameter types and arity ahead of any instruction.
Identical prompt code, opposite exposure. This is also why
Section~\ref{sec:validity} scrubs the disassembly symbol and not merely the
declaration: that is the channel the field actually leaks through.

Two limits on our own scope follow. First, we could not establish that the
practice is prevalent in \emph{modern-language} evaluations. We located no Dart,
Java, Kotlin, Swift, Go or Rust neural-decompilation evaluation other than our own
prior studies, so a prevalence claim we made in earlier work is withdrawn for want
of evidence rather than defended. The one non-C system verified, SCDBench, does
supply the original name with parameter types and arity as an explicit prompt
field, but its EVM target places selectors in the bytecode by design and it is a
single point. Second, what the survey does establish uniformly is the absence of
controls: \textbf{none of the six configurations reports an identifier-scrubbed or
signature-only arm}. The precondition this paper argues for is unmet everywhere we
looked, which is the general form of our contribution --- not that a particular
benchmark is broken, but that none of them can currently tell.

The evidence is uncomfortable in a specific way. When we remove the semantic
identifiers and leave everything else --- types, arity, the instruction stream,
the checkpoint, the compiler, the tests --- functional correctness collapses by a
factor of six \emph{while compilation improves}. The model does not become worse
at writing Dart. It becomes worse at writing the \emph{right} Dart. And in a
separate leakage-free study, an untrained baseline that sees only the signature
and never sees the binary at all is not measurably worse than any model we
trained on the binary.

\begin{table}[t]
\caption{Prompt channels in published neural decompilation evaluations. Each row
is grounded in released prompt-construction code, a released dataset, or a
verbatim prompt template. ``sig.'' distinguishes the \emph{original} source
signature (never supplied) from a decompiler-inferred one.}
\label{tab:survey}
\centering\scriptsize
\setlength{\tabcolsep}{3pt}
\begin{tabular}{@{}llllcc@{}}
\toprule
\textbf{system / split} & \textbf{name} & \textbf{params} & \textbf{sig.} & \textbf{ctrl} & \textbf{seeds}\\
\midrule
L4D-End / HumanEval-Dec.  & \texttt{func0}       & ---              & ---     & no & 1\\
L4D-Ref / HumanEval-Dec.  & \texttt{func0}       & \texttt{param\_N} & Ghidra  & no & 1\\
L4D / GitHub2025, ExeBench & \textbf{original}   & ---              & C++ types & no & 1\\
SALT4Dec. / Dec.-Eval, MBPP & \texttt{func}      & \texttt{var1..}  & ---     & no & 1\\
SK2Dec. / stripped         & \texttt{sub\_XXXX}  & \texttt{a1,a2}   & IDA     & no & 1\\
SCDBench / EVM-Solidity    & \textbf{original}   & ---              & \textbf{+types, arity} & no & 1\\
\midrule
\multicolumn{4}{@{}l}{\emph{totals}} & \textbf{0/6} & \textbf{0/4}\\
\bottomrule
\end{tabular}
\end{table}

\subsection{Contributions}
\begin{enumerate}[leftmargin=*,topsep=2pt,itemsep=1pt]
\item \textbf{A benchmark-validity result for neural decompilation}
(Section~\ref{sec:validity}): an inference-time ablation on a frozen checkpoint
shows that the semantic identifiers, not the binary, carry most of the measured
correctness, and an untrained signature-only baseline that never sees a binary is
not measurably worse than models fine-tuned on binaries.
\item \textbf{A measured resolution floor for \passatk}
(Section~\ref{sec:floor}). Re-running one frozen checkpoint under five sampling
seeds moves \passat{10} across $2$--$6$ of $174$ tasks ($\mathrm{SD}=1.79$). Most
differences reported in this literature --- including several of our own --- fall
below that floor, which we argue makes a seed-replicated floor a precondition for
interpreting any \passatk comparison.
\item \textbf{Aggregate reproducibility without compositional reproducibility}
(Section~\ref{sec:turnover}). Across every arm, \passatk counts are far more
stable than the task sets they summarise: our best arm's count varies by
$\mathrm{SD}=1.00$ while its five runs agree on only $5$ of the $25$ tasks they
collectively solve. Scoring the \emph{same} runs on compilation, which operates
at a much higher rate, reproduces at Jaccard $0.83$--$0.97$, which localises the
defect to the low-rate regime rather than to the harness or the benchmark. We
give a two-integer reporting fix.
\item \textbf{An inference-time channel decomposition}
(Section~\ref{sec:channel}). On a frozen checkpoint, supplying a typed contract moves
\compileat{10} from $126.0$ to $168.4$ of $174$ with \emph{no training}; removing
recovered constants collapses \passat{10} while leaving compilation statistically
unchanged; permuting which binary belongs to which task moves compilation and
not correctness; and withholding the binary entirely, while supplying the same
typed contract, maximises compilation while solving no more than baseline.
Supplying the contract at inference and fine-tuning on it both help; on the
primary endpoint the training step is the larger of the two, which contradicts a
claim we made in earlier work.
\item \textbf{Exploratory architecture and objective sweeps}
(Section~\ref{sec:arch}): sixteen controlled interventions with one significant
effect and a control-flow-edge effect whose sign reverses across seeds, together
with a rejection-sampling programme closed across six checkpoints and three
arrangements.
\item \textbf{A frontier ceiling on name-scrubbed data}
(Section~\ref{sec:frontier}), measured under a sealed, hash-pinned harness and
reported as a single-provider, single-date reference point.
\item \textbf{A documented evaluation-leakage post-mortem}
(Section~\ref{sec:leak}) and the protocol that replaced it
(Section~\ref{sec:protocol}).
\end{enumerate}

\subsection{Relationship to our prior work}
This paper builds on two earlier studies by the authors. A SANER ERA paper
introduced the Dart decompilation task and a first fine-tuned
pipeline~\cite{abualazm2026idiomatic}, and a journal study analysed fine-tuning
effectiveness and metric validity on that
pipeline~\cite{abualazm2026tosem}. Both report decompilation performance
\emph{using} a benchmark that supplies the function signature. The present paper
is the first to ask whether that benchmark is measuring decompilation at all,
and it answers in the negative. Section~\ref{sec:prior} summarises
results from~\cite{abualazm2026tosem} as context and claims no novelty for them;
Sections~\ref{sec:floor}, \ref{sec:validity}, \ref{sec:channel}, \ref{sec:arch},
\ref{sec:frontier} and~\ref{sec:leak} are new.

\section{Background and Related Work}
Classical decompilation descends from Cifuentes~\cite{cifuentes1994reverse}
through structural analysis~\cite{schwartz2013native}, control-flow
structuring~\cite{yakdan2015no} and platforms such as
BAP~\cite{brumley2011bap}. Neural approaches began with sequence
transduction~\cite{katz2018using,fu2019coda} and now centre on LLMs fine-tuned
on assembly--source pairs~\cite{llm4decompile}.

The current frontier adds structure or iteration atop a language model:
control-flow augmentation~\cite{liu2025cfadecllm}, hierarchical graph
abstraction~\cite{achamyeleh2026helios}, abstract logic
trees~\cite{wang2025salt4decompile}, skeleton-to-skin
generation~\cite{tan2025sk2decompile}, hierarchical attention over
assembly~\cite{jiang2025nova}, feedback-driven
refinement~\cite{liu2026autodecompiler}, and joint code and type
prediction~\cite{dramko2025idioms}. Our results engage this line directly and
mostly negatively: Section~\ref{sec:arch} reports a leakage-free null for
explicit control-flow and dataflow conditioning, and Section~\ref{sec:validity}
shows that the semantic identifiers --- which reach the model through the
channels catalogued in Table~\ref{tab:survey} --- account for most of the
measured success.

Evaluation has converged on \codebleu~\cite{ren2020codebleu}, \compileatk, and
\passatk~\cite{chen2021evaluating}, with benchmarks ported in the style of
MultiPL-E~\cite{cassano2023multiple}. Recent work questions whether such
benchmarks reflect real-world decompilation~\cite{decompilebench2025}; we supply
a mechanism for why they do not.

Dart is a demanding and under-studied target. Its AOT compiler performs
whole-program optimisation, erases names, and stores literals in a per-isolate
object pool addressed through a dedicated register, so constants do not appear
inline in the disassembly. Existing Dart tooling recovers structure but not
source~\cite{blutter,reflutter,jebdartaot}.

\section{Experimental Infrastructure}
\label{sec:infra}
This paper uses five corpora, and results are not comparable across them.
Table~\ref{tab:corpora} states which supports which claim.

\begin{table}[t]
\caption{Corpora. Counts are never compared across rows.}
\label{tab:corpora}
\centering\scriptsize
\begin{tabular}{@{}llrl@{}}
\toprule
\textbf{ID} & \textbf{origin} & \textbf{tasks} & \textbf{supports}\\
\midrule
HE-154 & HumanEval ported to Dart      & 154 & \S\ref{sec:validity}, \S\ref{sec:arch}\\
MF-175 & internal: LLM-gen $+$ legacy   & 175 & \S\ref{sec:floor}, \S\ref{sec:turnover},\\
       &                               &     & \S\ref{sec:channel}, \S\ref{sec:frontier}\\
CB-73  & Dart, held out                &  73 & \codebleu (prior work)\\
CP-126 & Dart, held out                & 126 & \compileatk (prior work)\\
FL-343 & Flutter ARM64 release APKs    & 343 & \S\ref{sec:arm64}\\
\bottomrule
\end{tabular}
\end{table}

\textbf{HE-154} is 154 of HumanEval's 164 problems with executable harnesses
(median 7 assertions). Two distinct defect populations affect it, and we separate
them because they have been conflated. First, 12 of the 154 supplied reference solutions fail their own harness. A
failing reference does not make its task unsolvable --- another implementation
may pass --- but it means those 12 tasks cannot be certified by their own gold
code, and we treat 142 as the reference-certified set. Three tasks fall in both
defect populations, so the strict union leaves 141. Second, and independently, 4 tasks carry hidden-test
defects that make them unscoreable under \emph{any} opaque contract: tests
binding an undefined wrapper (\#121, \#153), tests passing record arguments to a
\texttt{List<int>} contract (\#127), and an unintended string interpolation in
the test source (\#161). Those 4 are enumerated in a manifest frozen
independently of any treatment outcome; excluding them gives the
\textbf{valid-150} confirmatory set reported alongside the full 154 in
Section~\ref{sec:validity}. HE-154 supports the identifier ablation only.

\textbf{MF-175} is an \textbf{internally constructed, held-out x86-64 evaluation
suite}, not an external benchmark. It comprises 136 LLM-generated Dart tasks with
supplied test harnesses and 39 legacy standalone Dart SFT programs with generated
differential tests. Following signature neutralisation and a deterministic
semantic-component split, eligible tasks were individually compiled to
source-only AOT ELFs with Dart 3.12.2 and represented using their extracted
target-function F2. The surviving manifests identify $1{,}580$ training tasks and
$175$ evaluation tasks; a separate expanded file of $2{,}776$ \emph{training
rows} was derived from the training split and is not a task pool. MF-175 supports
the reseeding, stability, channel-decomposition and frontier analyses.

Two limitations follow, and both bear on how far its results generalise. The
retained records for the 39 legacy tasks do not identify their original
repositories or licences; we therefore describe them only as legacy internal
\texttt{dart\_all} tasks and make no project-level or real-world provenance
claim. More consequentially, \emph{a majority of the suite is model-generated}.
Synthetic Dart may lie closer to a decompiler's pretraining distribution than
human-written code, and its identifiers may be more uniformly descriptive than
those in production software --- either of which could inflate the
identifier-channel effects we measure. MF-175 therefore bounds what these
channels carry on this construction; it is not a sample of real-world
decompilation targets. Each task pairs a
reference implementation with an executable acceptance suite, and the references
pass their own harness $175/175$ --- the Rank-0 control of
Section~\ref{sec:floor}. One held-out task shares
byte-identical acceptance tests with a training row. That task is not inert: it compiles in every arm and every seed, and is solved
in every baseline, typed-contract and structure-permuted seed, while no
constants-stripped seed solves it. It is excluded from all confirmatory
analysis, and counts on MF-175 are reported on the leakage-clean
\textbf{174-task} subset. Its behaviour is itself a contamination datapoint: a
training row is recovered whenever the identifier and constant channels are
intact, and lost when the constants are removed. CB-73 and CP-126 are the
disjoint corpora used by the prior work summarised in Section~\ref{sec:prior}.

Functional comparisons are paired at the task level: McNemar's exact two-sided
test on discordant tasks, with 10{,}000-iteration bootstrap CIs. Across-seed arm
comparisons are \textbf{blocked, not independent}: every arm runs on the same
five seeds over the same tasks, so the valid randomisation is a paired sign-flip
over five seed blocks ($2^5=32$ arrangements), not a relabelling of ten free
observations. Its smallest attainable two-sided $p$ is $2/32=0.0625$, so
\textbf{no contrast in this study can reach $\alpha{=}0.05$ whatever its
magnitude}; six seeds would be the minimum for significance to be \emph{attainable}, not
guaranteed. We report the paired $p$ throughout
and treat $0.0625$ as the design's floor rather than as evidence. Rates carry Wilson intervals. \passatk for $k>1$ is
any-of-$k$ at the task level. \passat{1} on MF-175 is the success of the
first sampled candidate, not the $c/n$ estimator; it is reported for
completeness and carries no claim.

\textbf{What \compileatk measures, precisely.} We state this because several of
our results are otherwise easy to over-read. A candidate is concatenated with the
task's real test file and executed with \texttt{dart run}; a non-zero exit counts
as a compile failure \emph{only} when Dart emits a static front-end diagnostic,
while runtime exceptions and failed assertions count as compiled-but-not-passing.
Because the tests invoke the function by name with typed arguments and consume
its return value, and because Dart is statically typed, passing the front end
requires the emitted function's name, arity and parameter and return types to
conform to the call sites, and requires the body to type-check. \compileatk is
therefore an \emph{interface-conformance and type-consistency} metric. It is not
a test of whether the output is well-formed Dart, and it is not a proxy for
correctness. Any intervention that supplies the declared signature is supplying
the model with the answer to much of what this metric asks, and we read the
compilation column accordingly throughout.

\begin{table}[t]
\caption{Study matrix. Two distinct systems are used; no result crosses them.}
\label{tab:studies}
\centering\scriptsize
\setlength{\tabcolsep}{3pt}
\begin{tabular}{@{}llll@{}}
\toprule
 & \textbf{identifier (\S\ref{sec:validity})} & \textbf{channel/floor/stab.} & \textbf{frontier}\\
\midrule
corpus   & HE-154 / valid-150 & MF-174 & MF (full 175)\\
model    & Qwen3-8B$+$graph enc. & T5Gemma-2-4b-4b & DeepSeek-V4-Pro\\
weights  & \texttt{e8e872\ldots} & opt-348 \texttt{83d815\ldots} & provider API\\
seeds    & 42 (single) & 42--46 & one complete run\\
$k$ / samples & 10 & 10 & 10\\
Dart     & 3.11.5 & 3.12.2 & 3.12.2 (scoring)\\
prompt   & \texttt{antigravity-v2/v3} & sealed F2 contract & sealed F2 prompt\\
\bottomrule
\end{tabular}
\end{table}

The identifier study and the channel study run on \emph{different systems}
(Table~\ref{tab:studies}); no comparison in this paper crosses that boundary, and
agreement between them is convergence of two instruments, not replication on
one.

\section{Reseeding Variability Under This Protocol}
\label{sec:floor}
Every comparison in this paper, and every comparison in the literature it
critiques, is a difference between two \passatk counts on a benchmark of a few
hundred tasks. Before reporting any such difference we measured what this
protocol can resolve.

We re-ran a \emph{single frozen checkpoint} on the same 175 held-out tasks under
five sampling seeds, changing nothing else --- same weights, same task order,
same decoding parameters, same scorer, same evaluation file by SHA-256.

\begin{table}[t]
\caption{One checkpoint, the leakage-clean 174 tasks, five sampling seeds. Only the seed varies.}
\label{tab:floor}
\centering
\footnotesize
\begin{tabular}{@{}lrrrrrrr@{}}
\toprule
 & \multicolumn{5}{c}{seed} & & \\
\cmidrule(lr){2-6}
\textbf{metric} & 42 & 43 & 44 & 45 & 46 & \textbf{mean} & \textbf{SD} \\
\midrule
\passat{10}   & 6 & 6 & 6 & 4 & 2 & 4.8 & \textbf{1.79} \\
\compileat{10} & 123 & 124 & 134 & 125 & 124 & 126.0 & 4.53 \\
\passat{1}    & 2 & 1 & 2 & 1 & 2 & 1.6 & 0.55 \\
\bottomrule
\end{tabular}
\end{table}

\passat{10} ranges from $2$ to $6$ on an identical model
(Table~\ref{tab:floor}). A pre-registered interpretation gate --- fixed before
the seeds were drawn, and triggered by the observed spread --- declares
differences of this magnitude uninterpretable under this protocol. The practical
consequence is that \textbf{a difference of fewer than about three tasks cannot
be resolved here}, and we report every such difference as lying within the floor.

Two controls establish that this is sampling variance rather than a broken
instrument. First, scoring the reference implementations through the same
evaluator returns $175/175$ on compilation, \passat{1} and \passatk: the harness
registers a perfect result when one exists, so null outcomes are not an artifact
of measurement. Second, re-deriving the same arms under a later revision of the
extraction and scoring code reproduces every count exactly.

We report this as a contribution rather than a caveat, and we are careful about
whose results it constrains. The C-target literature surveyed in
Section~\ref{sec:scope} does not report \passatk at all: its metric is
re-executability, or RC/RE/TCP, on a \emph{single greedy decode}, so a decode-seed
floor does not transfer to those numbers directly. What does transfer is the
absence of any error bar --- none of the four papers reports seeds, repetitions,
variance or confidence intervals --- and the consequence surfaces in their
ablations, where component contributions of $0.1$--$0.3$\pp are reported and two
are sign-inverted, the ablated model beating the full one. Our floor bounds what
\passatk comparisons can resolve on corpora of this size, our own included; for
single-sample metrics the corresponding quantity is simply unreported.
Section~\ref{sec:protocol} makes a seed-replicated floor a reporting requirement.

\section{A Reproducible Number Over an Irreproducible Set}
\label{sec:turnover}
The floor of Section~\ref{sec:floor} is a statement about the \emph{count}. A
count can be stable while the set it summarises is not, and on this corpus it is.
For every arm of the channel study we recovered the identity of the solved tasks
in each of its five seeds and compared the sets pairwise
(Table~\ref{tab:turnover}).

\begin{table}[t]
\caption{Solved-set stability across five sampling seeds. Only the seed varies.
``all 5'' is the number of tasks solved in every seed, ``any'' the number solved
in at least one; Jaccard is the mean over all ten seed pairs. Counts reproduce
far better than the sets they summarise, and compilation reproduces far better
than correctness.}
\label{tab:turnover}
\centering
\scriptsize
\setlength{\tabcolsep}{3.5pt}
\begin{tabular}{@{}lrrrrrrr@{}}
\toprule
 & \multicolumn{5}{c}{\passat{10}} & \multicolumn{2}{c}{\compileat{10}} \\
\cmidrule(lr){2-6}\cmidrule(lr){7-8}
\textbf{arm} & \textbf{mean} & \textbf{all} & \textbf{any} & $J$ & $J$/chance & $J$ & $J$/chance\\
\midrule
$-$ constants      & 2.6  & 1 & 7  & 0.342 & 45.7 & 0.845 & 1.52\\
baseline           & 4.8  & 1 & 15 & 0.203 & 14.7 & 0.854 & 1.50\\
structure permuted & 5.2  & 1 & 12 & 0.302 & 20.4 & 0.830 & 1.62\\
$+$ typed contract & 10.0 & 4 & 24 & 0.375 & 13.1 & 0.973 & 1.04\\
\bottomrule
\end{tabular}
\end{table}

\textbf{The count is the stable part.} The typed-contract arm reports
$10.0\pm1.00$ solved tasks --- the tightest count in the study. Yet its five runs
solve $24$ distinct tasks between them and agree on $4$. A single run of that arm
therefore understates what the configuration can reach by more than half, and
overstates what it reaches \emph{reliably} by more than double, while reporting a
number that replicates to one task. The baseline is worse: two tasks are solved
in all five seeds, fifteen in at least one, and most of that union is solved by
exactly one run.

\textbf{The harness is not the cause.} The identical runs scored on compilation
put only $3\%$ of the compiling union in a single seed, against $56$--$71\%$ of
the solved union. Same generations, same seeds, same scorer, same tasks: whatever
scatters the solved set leaves the compiling set almost fixed, which excludes
nondeterminism in generation, in the compiler, or in evaluation order.

\textbf{But the raw Jaccard gap is not the measure of it.} Jaccard rises
mechanically with set size --- two sets covering $97\%$ of the tasks must overlap
heavily --- so the $0.30$ versus $0.97$ contrast is largely an artifact of the
rate, not evidence about it. Against the overlap expected for independent random
subsets of the same sizes, the ordering \emph{reverses}: solved sets run
$13$--$46{\times}$ chance while compiling sets run only $1.0$--$1.6{\times}$.
There is real task-level structure in what gets solved; the stable core is simply
small in absolute terms. The defensible statement is the absolute one --- of the
$24$ tasks the strongest arm ever solves, $4$ are solved every time --- and it is
the absolute figure that governs whether a task-level claim is safe.

This is the empirical form of the unmeasurable left tail: tasks whose per-sample
success probability is small relative to $1/k$ are individually indistinguishable
from unsolvable, and which of them surface is decided by the draw.

\textbf{What this costs the literature.} It does not threaten aggregate
comparisons: the between-arm gaps in Table~\ref{tab:channel} are several times
the count SD and survive. It does threaten every \emph{task-level} claim, and
those are common --- that a method fixes the harder cases, repairs a particular
failure mode, or unlocks a specific class of function. On this corpus, two runs
of one configuration disagree about roughly $70\%$ of what either solves. No such
claim is supportable at these rates without showing that the solved set itself
replicates.

The fix is cheap. Alongside the mean, report the number of tasks solved in
\emph{every} seed and in \emph{any} seed. Two integers make the gap visible, and
their ratio is a direct estimate of how much of a reported score is a stable
capability rather than a draw.

\section{The Benchmark Measures Its Own Identifiers}
\label{sec:validity}

\subsection{Design A: removing the semantic identifiers}
We freeze a checkpoint and hold the HE-154 tasks, the scoring stack and the
pinned compiler fixed, varying only the human-readable identifiers:
the original semantic HumanEval signature; \texttt{neutral\_exact}, which
preserves types and arity but replaces the function name with an opaque \fnz\ and
every parameter name with a positional letter; and \texttt{name\_only}, which
gives \fnz with no types.

We state the manipulation precisely, because it is broader than a function
rename and the breadth bounds what may be concluded. In all 154 tasks
\texttt{neutral\_exact} replaces the function name, all $200$ parameter-name
slots, the symbol in the disassembly headers and branch-target annotations, $15$
helper symbols, and the source path. What it leaves untouched is equally exact:
parameter types and arity are identical in $154/154$, and after masking symbol
annotations the instruction stream --- opcodes, operands, addresses --- is
byte-identical in $154/154$. The arm therefore removes \emph{every human-meaningful
identifier reachable by the model} while holding the binary evidence constant.
Effects below are attributable to that bundle, not to the function name alone.

\begin{table}[t]
\caption{Identifier ablation on a frozen checkpoint, seed 42. The function name,
all parameter names and the disassembly symbol change together; types, arity and
the instruction stream are held byte-identical.}
\label{tab:sig}
\centering\footnotesize
\begin{tabular}{@{}lrrrr@{}}
\toprule
\textbf{Signature shown} & \textbf{p@1} & \textbf{p@10} & \textbf{c@1} & \textbf{compiled}\\
\midrule
Semantic name      & 17.66 & \textbf{27.27} & 77.08 & 1{,}187/1{,}540\\
Typed, opaque \fnz & 0.71  & 4.55  & \textbf{84.61} & \textbf{1{,}303/1{,}540}\\
Name only          & 0.00  & 0.00  & 3.18  & 49/1{,}540\\
\bottomrule
\end{tabular}
\end{table}

Table~\ref{tab:sig} separates two contributions that the literature conflates,
including work predicting code and types jointly~\cite{dramko2025idioms}.

\textbf{Types and arity carry compilation.} Given an opaque typed contract the
model emits a function named \fnz\ with the requested arity in 1{,}537 of 1{,}540
samples, and 1{,}303 of those compile against the tests --- \emph{better} than the
named comparator by $7.5$\pp. Syntactic competence does not degrade.

\textbf{The identifiers carry functional correctness.} Removing them
collapses \passat{10} from $27.27\%$ to $4.55\%$. Of the comparator's 42 solved
tasks, \texttt{neutral\_exact} retains 7 --- zero gained, 35 lost. Roughly $83\%$
of ``solved'' tasks were solvable only with the semantic identifiers present.
That the floor is
reachable rules out a broken harness: 151 of 154 neutral references compile and
142 pass. The scrubbed generations are not broken Dart; they are
\emph{semantically wrong} Dart.

\textbf{The result is invariant to the benchmark's defects.} Repeating the
comparison on the valid-150 set of Section~\ref{sec:infra}, which drops the four
contract-defective tasks, gives $28.0\%$ versus $4.67\%$ --- and \emph{the same}
$35$ losses, $0$ gains and $7$ retained. The excluded tasks all sit among the
$112$ that neither arm solves, so removing them moves the denominator and nothing
else. The reference-certified 142 cannot be recomputed without per-task
rescoring, but it is bounded: even if all 12 reference-failing tasks lay among
the 35 losses, the containment would remain at least 23 losses to 0 gains. We
report all denominators for this reason: the effect does not depend on which
validity criterion a reader prefers.

We locate this figure against the standard of Section~\ref{sec:turnover} rather
than leave it to be questioned. The $83\%$ is a ratio of two counts, $42$ and
$7$, and the count is the quantity Section~\ref{sec:turnover} finds stable; the
collapse is a factor of six, far outside any floor. The containment it rests on
--- $35$ losses and \emph{zero} gains --- is a task-level statement, but it is
the opposite of the fragile case: the discordance is unanimous across all $35$
pairs (exact two-sided $p{\approx}6\times10^{-11}$), at a $27.27\%$ success rate
where Table~\ref{tab:turnover} shows solved sets to be far more stable than in
the low-rate arms. The constraint of Section~\ref{sec:turnover} applies to
task-level claims resting on few discordant pairs at low success rates, which
this is not. What does remain conditional on this seed is the \emph{identity} of
the seven retained tasks, and we make no claim about which they are.

\subsection{Design B: an untrained signature-only baseline}
Design A varies a prompt on a fixed model. The complementary question is whether
\emph{training} on binaries buys anything over a policy that never sees one. In
the leakage-free study of Section~\ref{sec:arch} we compared three
binary-consuming arms against a signature-only baseline that is the
\textbf{untrained base model}: no fine-tuning, no adapter, prompted with the
signature alone. It reaches \passat{10} of $0.234$ ($36/154$).

\begin{table}[t]
\caption{Every arm against a baseline that sees only the signature and never the
binary. Paired task bootstrap; McNemar exact.}
\label{tab:sigonly}
\centering\footnotesize
\begin{tabular}{@{}lrlr@{}}
\toprule
\textbf{Arm vs.\ signature-only} & \textbf{$\Delta$p@10} & \textbf{95\% CI} & \textbf{$p$}\\
\midrule
Trained on raw assembly & $-0.0195$ & $[-0.065,\,+0.026]$ & 0.549\\
Graph prefix, no edges  & $+0.0195$ & $[-0.020,\,+0.058]$ & 0.549\\
Graph prefix + CFG      & $+0.0260$ & $[-0.013,\,+0.071]$ & 0.388\\
\bottomrule
\end{tabular}
\end{table}

\textbf{No arm is distinguishable from the signature-only baseline.} A model
fine-tuned on raw assembly is nominally \emph{worse} than an untrained model that
never sees the binary; on \passat{1} that deficit's interval excludes zero. The
graph-conditioned arms are nominally better by two to three points, with
intervals crossing zero.

We are explicit about what this comparison is and is not. The baseline differs
from the other arms in \emph{two} respects --- it is untrained and it is shown a
different input --- so it does not isolate the binary channel and cannot support
a decomposition. What it does support is a bound, and the bound is the
uncomfortable part: several hundred GPU-hours of fine-tuning on binaries did not
produce a policy measurably better than the base model reading a signature.
Designs A and B are methodologically independent, and both are consistent with
the benchmark being insensitive to binary evidence; neither establishes that the
binary is uninformative, and we do not claim it.

\section{Exploratory: Architecture and Objective Interventions}
\label{sec:arch}
A natural response is that the architecture was inadequate. We ran a
pre-registered sweep with freeze rules fixed in advance and a
reviewer-objection matrix mapping each plausible criticism to a control, rather
than searching until a configuration won.

A three-seed core ablation (prefix with no edges, $+$CFG, $+$CFG$+$DFG; seeds
42--44) puts the no-edge arm at the best mean \passat{10} ($.260$) with every
edge effect's hierarchical-bootstrap interval crossing zero. The instructive part
is the instability: at seed 44, CFG versus no-edge is 1 gain / 8 losses (exact
$p{=}0.0391$) while adding DFG to CFG is 7 gains / 0 losses ($p{=}0.0156$) --- two
nominally significant effects in opposite directions within one seed. A direction
that reverses across seeds is not evidence of a null; it is no evidence of a
reproducible effect either way.

Table~\ref{tab:sweep} reports the component isolations. Encoder identity,
message-passing depth, prefix width, hierarchical regions, global attention and
positional encodings all leave \passat{10} statistically unchanged.

\begin{table}[t]
\caption{Architectural interventions (seed 42, paired task bootstrap). Exactly
one is significant, and it removes the channel rather than restructuring it.}
\label{tab:sweep}
\centering\footnotesize
\begin{tabular}{@{}lrr@{}}
\toprule
\textbf{Intervention} & \textbf{$\Delta$p@10} & \textbf{$p$}\\
\midrule
No-GINE vs.\ no-edge GINE            & $+.0065$ & 1.000\\
No-GINE vs.\ CFG+DFG GINE            & $+.0260$ & 0.219\\
CLAP-ASM vs.\ GraphCodeBERT encoder  & $-.0065$ & 1.000\\
Four block vectors vs.\ one CLS      & $+.0000$ & 1.000\\
Eight block vectors vs.\ one CLS     & $-.0130$ & 0.688\\
Hierarchical regions vs.\ block prefix & $-.0195$ & 0.375\\
Region maximum 16 vs.\ no regions    & $+.0130$ & 0.688\\
No global attention                  & $-.0130$ & 0.625\\
No block positions                   & $-.0195$ & 0.375\\
Two- vs.\ four-layer GINE            & $+.0065$ & 1.000\\
Cross-task permuted graph            & $-.0130$ & 0.500\\
Shuffled block order                 & $-.0065$ & 1.000\\
\midrule
\textbf{Zero the prefix gate}        & $\mathbf{-.0455}$ & $\mathbf{0.016}$\\
\bottomrule
\multicolumn{3}{@{}p{0.92\linewidth}@{}}{\footnotesize Three further contrasts
from the same sealed sweep are omitted for space; all are null
($-.0065$, $p{=}1.000$; $+.0000$, $p{=}1.000$; $-.0195$, $p{=}0.508$). Including
them, 16 contrasts were run over 15 intervention arms, and the single nominally
significant result below is the only one among all 16.}
\end{tabular}
\end{table}

\textbf{Multiplicity.} We report the correction rather than wait to be asked.
With 16 contrasts, a Holm threshold at the smallest rank is ${\approx}0.003$, and
the prefix-gate effect at $p{=}0.016$ does not survive it; under the global null,
seeing at least one $p<0.05$ among 16 tests has probability ${\approx}0.56$. We
retain the arm not as an independently significant finding but because it was
predicted in advance, its discordance is unanimous ($0$ gains, $7$ losses), and
it converges with Designs A and B. Read alone, it is consistent with chance.

\textbf{The one nominally significant cell does not survive correction.} Zeroing the prefix
gate --- deleting the binary channel outright --- costs $0.0455$ with 0 gains and
7 losses ($p{=}0.0156$), and freezing the encoder drops \compileat{1} from
$0.78$ to $0.48$. Meanwhile substituting \emph{another task's} graph costs only
$0.0130$ ($p{=}0.50$) and shuffling block order costs $0.0065$ ($p{=}1.00$). The
decoder uses the channel, but is largely indifferent to what it contains: the
learned assembly-to-prefix compression is load-bearing, the explicit topology is
not.

\subsection{Disclosure}
This sweep is exploratory, and we treat it as such: the core ablation has three
of five planned seeds, the component isolations are single-seed diagnostics, and
the study's own provenance audit header reads \textsc{fail} across its 38 runs
(10 showing an aggregate/candidate-replay disagreement, resolved by treating
candidate-level replay as canonical). None of its conclusions is confirmatory;
the certified result of this paper is the pre-registered channel test of
Section~\ref{sec:seed47}, not anything in this section.

\section{Prior Findings on This Corpus}
\label{sec:prior}
We cite~\cite{abualazm2026tosem} as context only, claim no novelty for it, and
--- because it remains under journal review --- do not reproduce its results
here. In summary, it reports fine-tuning configurations without
functional-correctness gains and a divergence between surface metrics and
correctness on this corpus. Those findings establish that training on this
corpus does not improve functional correctness and that surface metrics do not
track it. They do not address a prior question: what the benchmark measures when
it does report a score. That question is this paper's subject.

\section{A Frontier Reference Point on Identifier-Scrubbed Data}
\label{sec:frontier}
Our first frontier measurement reported $5/60$. An internal audit ruled it
indefensible --- assembly truncated at 12{,}000 characters so 33 of 60 prompts
lost the function tail, anonymous closures omitted by the symbol regex, a cache
that was not content-addressed, and unknown effective $k$ with silent failures.
We never report it as a result.

\begin{sloppypar}
The replacement harness pins a paired manifest over 175 ordered tasks and
enforces eight invariants, all verified:
{\footnotesize\texttt{provider\_prompts\_exclude\_gold\_and\_tests}},
{\footnotesize\texttt{no\_prompt\_truncation}},
{\footnotesize\texttt{measure\_only\_seals\_verified}},
{\footnotesize\texttt{prompt\_artifact\_hashes\_verified}},
{\footnotesize\texttt{same\_acceptance\_tests}},
{\footnotesize\texttt{same\_f2\_grammar}},
{\footnotesize\texttt{same\_ordered\_175\_tasks}},
{\footnotesize\texttt{same\_prompt\_budget\_contract}}. Evaluator, Dart binary, tokenizer and
prompt artifacts are SHA-256 pinned, each run carries a 256-bit completion
attestation, and the harness refuses to emit a score for an incomplete run.
\end{sloppypar}

On the complete run --- 175/175 tasks, 1{,}750/1{,}750 slots --- DeepSeek-V4-Pro
at $k{=}10$ solves \textbf{4 tasks}, \passat{10} $2.29\%$. We separate the two
denominators explicitly, because conflating them is easy and we have done it:
at the \emph{task} level \compileat{10} is $86/175 = 49.1\%$, while at the
\emph{candidate} level $300/1{,}750 = 17.1\%$ of generations compile and
$11/1{,}750 = 0.63\%$ pass. Every slot carries \texttt{finish\_reason: stop}, so
the score is not a truncation artifact. Our 8B student reaches \compileat{10} of
$69.1\%$ on the same tasks: whatever the frontier model contributes here, it is
not syntactic competence. This is one provider's model on one date under one
decoding configuration; it is a reference point, not a ceiling. No score exists
for the other frontier models we queried --- they reached preflight or quota only,
and absence of a score is not a score of zero.

\section{The Information Channel}
\label{sec:channel}
Sections~\ref{sec:validity} and~\ref{sec:arch} establish that the name carries
the score and that restructuring the model does not. This section asks the
constructive question: \emph{which} channel carries recoverable signal. We
decompose the input on a frozen checkpoint, varying one channel at a time.

One chronological fact is disclosed rather than left in the artifact: seed 42
predates the other four --- its baseline is the sealed epoch-ablation run and it
served as diagnostic context while the interventions were designed --- whereas
seeds 43--46 are fresh replications. Excluding seed 42 changes no direction in
any contrast. Each intervention was replicated across five sampling seeds (42--46); the typed
contract, baseline and structure-permutation arms carry a sixth (47) added under
the pre-registered amendment of Section~\ref{sec:seed47}. Table~\ref{tab:channel}
gives the decomposition; every effect is tested by an exact paired sign-flip over
the seed blocks, the valid randomisation given that all arms share the same
seeds.

\textbf{What the six-seed test establishes.} With five blocks the paired floor
was $2/2^5=0.0625$, so no contrast could reach $\alpha{=}0.05$ whatever its
magnitude. We therefore pre-registered one sixth seed with a \emph{single}
confirmatory test --- the typed contract against baseline on \passat{10} --- and a
hard stop (Section~\ref{sec:seed47}). Seed 47's difference is $+6$ tasks, the
sixth positive block of six, giving $p{=}2/2^6=0.031$: the typed-contract effect
on functional correctness is significant. Its compilation effect is
sign-consistent over the same six blocks at the same $p$, reported as a secondary
descriptive rather than a second confirmatory test. Structure permutation is
sign-consistent on compilation across six blocks ($-8.7$ tasks) and null on
correctness; removing constants stays at $-2.2$ tasks ($p{=}0.25$), directional
and below resolution. The paper's own standard --- replicate before interpreting
--- is now met for the one effect it certifies.

\begin{table}[t]
\caption{Input-channel decomposition on a frozen checkpoint, leakage-clean 174
tasks, six sampling seeds for the confirmatory arms and five elsewhere, as
marked. Each channel has a distinct signature, and two of the six tested cells
show no detected effect.}
\label{tab:channel}
\centering
\footnotesize
\begin{tabular}{@{}lcc@{}}
\toprule
\textbf{input} & \textbf{\passat{10}} & \textbf{\compileat{10}} \\
\midrule
baseline                 & $5.0 \pm 1.67$  & $125.8 \pm 4.07$ \\
$+$ typed contract       & $\mathbf{10.3 \pm 1.21}$ & $168.0 \pm 2.19$ \\
$-$ constants$^\ddagger$ & $\mathbf{2.6 \pm 1.14}$  & $124.2 \pm 4.66$ \\
structure permuted       & $5.0 \pm 1.79$  & $\mathbf{117.2 \pm 2.93}$ \\
typed contract, no binary$^\dagger$ & $4.0$ \tiny{(5,\,3)} & $\mathbf{173.5}$ \tiny{(174,\,173)} \\
\midrule
\multicolumn{3}{@{}l@{}}{\emph{effect vs.\ baseline (exact paired sign-flip, two-sided;}}\\
\multicolumn{3}{@{}l@{}}{\emph{six-seed floor $2/2^6=0.031$)}}\\
$+$ typed contract  & $+5.3$, $p{=}\mathbf{0.031}$ & $+42.2$, $p{=}0.031$ \\
$-$ constants$^\ddagger$ & $-2.2$, $p{=}0.25$   & $-1.8$, $p{=}0.50$ \\
structure permuted  & $0.0$, $p{=}1.00$    & $-8.7$, $p{=}0.031$ \\
\bottomrule
\multicolumn{3}{@{}p{0.94\linewidth}@{}}{\footnotesize Baseline, typed contract
and structure permutation are six seeds (42--47); $^\ddagger$constants is five
(42--46), its status unchanged by a sixth block. $^\dagger$Two seeds only,
per-seed counts in parentheses: a gold-derived oracle control with an
out-of-distribution empty input, reported without a permutation effect and
without any equivalence claim (Section~\ref{sec:nobinary}).}
\end{tabular}
\end{table}

\subsection{Types, supplied at inference, without training}
Replacing the input's signature with a typed but opaque contract --- \fnz\ with
declared return and parameter types, parameter names reduced to
\texttt{p0},\,\texttt{p1},\,\ldots\ --- and changing nothing else raises
\compileat{10} from $126.0\pm4.53$ to $168.4\pm2.19$ tasks of 174
($+42.4$) and \passat{10} from $4.8\pm1.79$ to $10.0\pm1.00$
($+5.2$); both reach complete separation across the five seeds, so both return
the paired floor of $0.0625$. On the paired seed-42 comparison the compilation gain
is $49$ tasks against $3$ losses; at candidate level compilation rises from
$578/1{,}740$ to $1{,}323/1{,}740$, and diagnostics attributable to wrong arity
or wrong types fall from $1{,}081$ to $312$.

We note that the functional gain is visible only once seeds are replicated. At a
single seed the difference is $6\rightarrow9$ tasks, inside the floor of
Section~\ref{sec:floor} and not interpretable; across five seeds it is a
$5.2$-task effect at the paired floor of $p{=}0.0625$. This is the concrete case for the
reporting requirement of Section~\ref{sec:protocol}.

The comparison that matters is what training adds on top, and here we correct a
claim we have made in earlier work. A model \emph{fine-tuned} on the same typed
contract reaches \compileat{10} of $171$ against $170$ --- one task --- which we
previously reported as showing that essentially the whole attainable gain is
information rather than optimisation. That reading selected the secondary metric.
On the primary endpoint at the same seed, \passat{10} goes from $10$ to $14$. The
full triple in functional terms is $6\rightarrow9\rightarrow14$: supplying the
contract adds three tasks, training on it adds five. Neither step is individually
resolvable (exact McNemar $p{=}0.27$ for the training step), and the honest
summary is that information and optimisation contribute comparably here, with the
data unable to separate them.

\subsection{A three-way dissociation}
The remaining two interventions do not simply produce smaller versions of the
same effect. Each moves a different metric, and each leaves the other unchanged.

\textbf{Constants: correctness moves; no compilation effect detected.} Deleting recovered string and numeric
literals, with the structural payload byte-identical, reduces \passat{10} from
$4.8\pm1.79$ to $2.6\pm1.14$ ($-2.2$, $p=0.25$) while \compileat{10} shows no
detected difference at $124.2\pm4.66$ ($-1.8$, $p=0.50$). On the 96 tasks
that actually carry literals the seed-42 count falls from 4 solved to 0. As a
control, the 79 tasks containing no literals return predictions
\emph{byte-identical} to baseline, confirming that the intervention is surgical
and the harness deterministic. Stripped of constants the model still writes
Dart that compiles; it stops writing Dart that is correct.

\textbf{Structure: compilation moves; no correctness effect detected.} An exact derangement, in which each
task receives another task's assembly and control-flow graph while keeping its
own constants, leaves \passat{10} statistically indistinguishable from baseline
($5.2\pm1.92$ versus $4.8\pm1.79$; $+0.4$, $p=0.875$) while reducing
\compileat{10} to $117.4\pm3.21$ ($-8.6$, $p=0.0625$). Running the model on an
entirely different function's binary changes functional correctness by
$+0.4$ tasks --- a nominal \emph{gain}, well inside the floor.

We note what this arm does \emph{not} show. The derangement preserves each
task's own recovered constants and replaces only the structural payload, so it
isolates the identity of the control-flow evidence rather than removing the
binary channel. Taken with the constants ablation above, the two arms together
explain each other: permuting structure is nearly free \emph{because} the
correctness-carrying channel is left intact.

We label this an \emph{exploratory channel pattern}, not a dissociation. Across
five paired seeds the three contrasts point in different directions on the two
metrics, but the design establishes no contrast at $\alpha{=}0.05$ and
non-significance in the remaining cells is not equivalence. The defensible
summary is that the typed contract moves both metrics by large margins, removing
constants moves correctness by $-2.2$ tasks with no compilation effect detected,
and permuting structure moves compilation by $-8.6$ tasks with no correctness
effect detected --- all of it directional evidence awaiting a six-seed test.

Across all four arms and every seed, the mean number of distinct programs emitted
per task lies between $9.98$ and $10.00$ of 10, so none of these effects is
attributable to degenerate decoding.

\subsection{The pre-registered sixth seed}
\label{sec:seed47}
The five-seed design cannot certify any effect: with all arms on the same seeds,
the valid randomisation is a paired sign-flip over the blocks, and at five blocks
its smallest two-sided $p$ is $2/2^5=0.0625$. Rather than report the programme's
central channel result as ``directional'' indefinitely, or add seeds until
something crossed a threshold, we sealed an amendment before generating a single
new token. It fixes \emph{one} confirmatory test --- the typed contract against
baseline on \passat{10}, the paper's declared primary endpoint --- one fresh seed
(47), and a hard stop: no seed beyond 47, whatever the result. Constants was
excluded by that amendment on identifiability grounds, its five-seed difference
vector containing zero blocks that no sixth block can make significant. The
amendment also records the two operational deviations forced by a smaller GPU ---
a generation batch reduced from ten to five, applied uniformly, and a
corresponding fix to a chunked-sampling defect in the harness --- both fixed
before any seed-47 output was scored, so seed 47 is one well-defined block rather
than a moving target. The regenerated tokenizer was accepted only after hashing
to the sealed identity byte-for-byte.

Seed 47 returned a typed$-$baseline \passat{10} difference of $+6$ tasks, the
sixth positive block of six ($3,5,4,5,9,6$); the exact paired test gives
$p{=}2/2^6=0.031$. The result is reported exactly as registered, and the two
compilation contrasts that are also sign-consistent over six blocks are reported
beside it as secondary descriptives, not as further confirmatory tests.

\subsection{Withholding the binary entirely}
\label{sec:nobinary}
The three arms above perturb the binary. The limiting case removes it. We
pre-registered a control --- hypotheses, thresholds, contrasts and input hashes
sealed before any generation --- that preserves the typed contract byte-for-byte
and replaces the entire binary payload with the empty string, task-invariant, at
the same frozen checkpoint.

We predicted $\compileat{10}\in[160,172]$ and $\passat{10}\in[1,3]$. Both
predictions were wrong, in the same direction: compilation came in at $174$ and
$173$ of the leakage-clean $174$ (the prediction was registered on the full
175), and correctness at $5$ and $3$. The registered branch that fired
is \emph{substantial type- or benchmark-prior performance}, which obliges us to
re-interpret the typed-contract result rather than report the control as a floor.

\begin{table}[t]
\caption{Candidate-level rates, leakage-clean 1{,}740 candidates per arm, seed
42. Conformance saturates while correctness stays at baseline level.}
\label{tab:candlevel}
\centering\footnotesize
\begin{tabular}{@{}lrr@{}}
\toprule
\textbf{input} & \textbf{compiled} & \textbf{passed}\\
\midrule
baseline                     & 33.2\% & 0.52\%\\
$+$ typed contract           & 76.0\% & \textbf{1.15\%}\\
typed contract, no binary    & \textbf{89.0\%} & 0.57\%\\
\bottomrule
\end{tabular}
\end{table}

\textbf{Conformance saturates; correctness does not move.}
Table~\ref{tab:candlevel} shows it on a denominator ten times larger than the
task-level one: deleting the binary and keeping the contract gives the highest
compilation rate we measured while candidate-level correctness stays at baseline
level ($0.57\%$ versus $0.52\%$) --- an earlier draft called this arm the least
correct, which the leakage-clean denominator does not support, and we withdraw
that reading. Given what \compileatk actually asks
(Section~\ref{sec:infra}), the compilation half is close to definitional --- the
model was handed the declared signature, which is most of what the front end
checks --- and we claim nothing more from it than that. The informative half is
that near-perfect interface conformance buys \emph{nothing} in correctness:
candidate-level correctness is statistically indistinguishable from baseline, and
at task level this arm solves no more than the baseline that reads the binary
without types, and fewer than half as many as the arm given both.

\textbf{The outputs are wrong, not malformed.} Mean distinct candidates per task
are $9.95$ and $9.97$ of $10$, so this is not decoding collapse. Every frequent
failure diagnostic is a runtime assertion --- mismatched list, map or set
lengths, \texttt{false != true}, \texttt{RangeError} --- rather than a static
error. The programs compile against the real tests, execute, return values of the
declared type, and compute the wrong thing.

\textbf{What the control does not license.} Against the same-seed typed$+$F2 arm,
the correctness deficit is directionally consistent in both seeds and significant
in one ($3$ versus $8$ discordant, $p{=}0.227$; $1$ versus $10$, $p{=}0.012$);
our pre-registration required significance in both, so we do not claim a
quantified F2 correctness contribution from this control. With types held
constant, removing the binary gains $5$ compiling tasks and loses $0$ and $1$
respectively ($p{=}0.0625$, which is the \emph{minimum attainable} $p$ at five
discordant pairs, and $p{=}0.219$), so the compilation direction replicates but
neither seed resolves it. The arm is a gold-derived oracle: it assumes perfect
type and arity recovery while evaluating no recovery front end, and the empty
input is out of distribution for this policy. We make no equivalence claim
against baseline in either metric, and the $4.0$ correctness figure is
baseline-\emph{scale}, not baseline-equivalent --- the two arms' solved sets
overlap in one task.

\subsection{One benchmark, four systems}
Table~\ref{tab:unified} places every system we measured on the same 175
name-scrubbed tasks. Every row was re-derived from the raw generation artifacts
under two toolchain revisions, returning identical counts.

\begin{table}[t]
\caption{Descriptive cross-system reference points, full 175-task runs predating
the single-task exclusion, $k{=}10$. The ordering is consistent with the channel
experiments; the table itself, which mixes model, training, input and scale,
supports no causal claim.}
\label{tab:unified}
\centering
\footnotesize
\begin{tabular}{@{}llrr@{}}
\toprule
\textbf{system} & \textbf{input} & \textbf{\passat{10}} & \textbf{\compileat{10}} \\
\midrule
Decoder-only 8B    & assembly only        & 1  & 103 \\
Decoder-only 8B    & +\,recovered consts. & 4  & 107 \\
Frontier reasoning & +\,recovered consts. & 4  & 86  \\
Decoder-only 8B    & +\,oracle consts.    & 10 & 121 \\
Encoder--decoder   & assembly + CFG       & 7  & 124 \\
Encoder--decoder   & +\,typed contract    & 14 & 171 \\
\bottomrule
\end{tabular}
\end{table}

A frontier reasoning model and a fine-tuned 8B decoder given the same recovered
constants both solve four tasks. We present that as a coincidence of magnitude
rather than an established equivalence: the solved sets were not compared, both
intervals are wide, and Section~\ref{sec:turnover} shows that at these rates two
runs of a \emph{single} system already disagree about most of what either solves,
so equal counts are especially weak evidence of equal behaviour. The defensible reading of the table is the column
ordering --- \passat{10} moves with what the input contains, and barely with what
consumes it.

\subsection{Objective-side interventions do not substitute}
If the deficit were a training-signal problem, objective-side work should close
it. Reinforcement learning never reached a step (seven launches, hard OOM at
${\sim}94$\,GiB). Dense knowledge distillation is not constructible from a
top-$k$ API. Expanded synthetic SFT changed \passat{10} by $+0.0065$
$[-0.020,+0.039]$.

Two further programmes are often described as null results and should not be:
reinforcement learning never completed a training step (seven launches, hard OOM
at ${\sim}94$\,GiB) and dense distillation was not constructible from a top-$k$
API. Those are feasibility failures, not measurements, and we withdraw them from
the paper's negative findings.

Rejection-sampled fine-tuning we pursued to exhaustion. Across six checkpoints,
three arrangements (a single pass, sequentially stacked passes, and one folded
pass over the union of all harvested targets) and two replay conditions,
\passat{10} ranged from $12$ to $18$ against a baseline of $14$ --- every value
inside the floor. One arrangement did produce a measurable effect, in the wrong
direction: stacking a second pass on the first left \compileat{10} and
\passat{1} unchanged or higher while \passat{10} fell from $18$ to $12$, with
mean distinct candidates per task falling from $10.00$ to $9.45$. The policy
became more reliable on what it already produced and lost the tail draws that
\passatk depends on. Folding the same data into a single pass from the
supervised checkpoint removed the collapse ($9.98$) and produced no gain.

Against all of this, the information-side lever works. Prepending constants
recovered from the binary raises \passat{10} from $0.57\%$ to $5.71\%$ with
\compileat{10} at $69.1\%$. We state it carefully: that arm reconstructs
constants by regex over \emph{gold source} and is an oracle upper bound
establishing the bottleneck is \emph{relievable}, not a deployable system. Our
binary-grounded extractor reaches $2.29\%$, succeeding on $1{,}746$ of $1{,}755$
tasks at string recall $0.465$ micro against its own $\geq0.8$ acceptance bar.

\subsection{Cross-ISA observation, and two reasons not to lean on it}
\label{sec:arm64}
On 343 held-out functions from real obfuscated Flutter release APKs targeting
ARM64 --- production binaries rather than a HumanEval port --- \passat{10} is
$5.54\%$ (19 tasks) and \compileat{5} is $93.07\%$. Pre-registered length bins
fall monotonically to zero, and of 168 tasks with ${\geq}200$ instructions none
was solved, while \compileat{1} decays only gently ($0.65\rightarrow0.39$). The
model keeps producing compilable code and stops producing correct code, which is
the same qualitative pattern the x86 decomposition resolves into channels.

We give two reasons, both drawn from this paper's own standards, for reporting
that number as an observation rather than a result. First, the corpus retains
original function names and the run supplied the full semantic signature, so
$5.54\%$ is measured under exactly the condition Section~\ref{sec:validity}
criticises. Second, it is a \emph{single seed}. By Section~\ref{sec:turnover} a
lone draw at a $5.5\%$ success rate fixes neither the count nor the identity of
the 19 solved tasks, and we have not replicated it. Either defect alone would
disqualify the figure as evidence about the ISA; together they make it a
length-scaling observation and nothing more.

We state what the replication now requires, because it is unusually cheap and the
instrument is unusually clean. The frozen checkpoint is recoverable and the four
signature views are built and validated ($343/343$ rows, zero residual names,
assembly byte-identical to the unmodified arm), and the signature mode and a
binary-channel ablation are both existing configuration flags, so the study is
inference-only with no code change. More important than cost, that architecture
separates the two channels \emph{physically}: the signature is tokenised text
while the binary enters as a learned soft prefix of embeddings, so unlike the
single serialised prompt used in Sections~\ref{sec:validity}
and~\ref{sec:channel}, leakage between the channels is structurally impossible
and zeroing the binary is exact.

One limit is worth stating in advance. Baseline compilation on that corpus is
$339/343$, so \compileatk has four tasks of upward headroom and interventions
that \emph{raise} compilation --- including the effect of
Section~\ref{sec:nobinary} --- cannot be reproduced there. The correctness half
of the decomposition is testable; the compilation half is testable only
downward.

\section{A Post-Mortem on Evaluation Leakage}
\label{sec:leak}
An internal audit found that our historical prompt builder extracted up to ten
assertion lines from each benchmark row's \texttt{tests} field into the policy
prompt; the evaluator then appended the same field to the generated code and
executed it. All 154 rows leaked, 38 saturated the ten-line cap, and 133 of 154
tasks (86\%) had their entire suite inside the window. Truncation preserved the
fixed prefix and trimmed \emph{assembly}, so the leak survived precisely where
the binary evidence did not.

Every archived functional pool affected has been withdrawn. Because no historical
run summary pinned a prompt hash, source commit or dataset hash, and file
timestamps are not proof of provenance, the exclusion was applied as a blanket
rule --- covering pools that may well be clean --- with the affected tables
retained as audit records rather than as results. The corpora underlying the
prior studies of Section~\ref{sec:prior} carry no scoring-test field and are not
affected by \emph{this specific} leak; that phrasing is deliberate, as it is not
a clean bill of health. All results in
Sections~\ref{sec:validity} and~\ref{sec:arch} postdate the fix and run under
the schema \texttt{antigravity-v2-no-test-hints} with
\texttt{rows\_leaking\_scoring\_tests: 0}.

\section{Recommendations}
\label{sec:protocol}
\begin{enumerate}[leftmargin=*,topsep=2pt,itemsep=1pt]
\item \textbf{Report an identifier-scrubbed arm.} Scrub the function name, the
parameter names \emph{and the symbol in the disassembly}; a benchmark that
supplies any of them cannot separate decompilation from reading a specification.
On our data the two differ by a factor of six, and the disassembly symbol is the
channel through which published C evaluations actually leak it
(Table~\ref{tab:survey}). None of the six configurations we surveyed reports such
an arm. This applies with particular force to HumanEval ports, whose reference
solutions are near-certainly in pretraining data.
\item \textbf{Include a signature-only baseline that never sees the binary.} It
costs \emph{no} training run --- the untrained base model suffices, and in our
study it was not beaten --- and it is the cheapest way to show the binary
contributes.
\item \textbf{Report the sign stability of effects across seeds}, not only
confidence intervals. A reversing sign and a wide interval are different
findings.
\item \textbf{Report tasks solved in \emph{every} seed and in \emph{any} seed
beside the mean.} Two integers, and they expose the gap between a reproducible
score and a reproducible capability (Section~\ref{sec:turnover}). Withhold any
claim about \emph{which} tasks a method fixes until that gap is shown to be
small.
\item \textbf{Decompose best-of-$k$ surface metrics.} A rise in a max-over-$k$
score computed only on compiling candidates may be a coverage effect; report the
compiled-only mean beside it.
\item \textbf{Fail closed on test content in prompts}, and record a
prompt-schema string in provenance.
\item \textbf{Pin artifacts and refuse partial scores.} A budget-truncated
denominator is not a random sample.
\end{enumerate}

\section{Threats to Validity}
\textbf{Construct.} That the benchmark measures its own identifiers rather than
decompilation is the paper's subject rather than a caveat to it. Two limits on
that claim are ours to state. First, the ablation removes a \emph{bundle} ---
function name, parameter names, disassembly symbol --- so it licenses no
attribution to any one of them; the parameter names alone are strong task cues on
HumanEval. Second, a descriptive identifier is simultaneously a retrieval key and
a specification, and no arm here separates memorisation from specification-driven
inference. The discriminating designs are paraphrased and misleading identifiers,
and tasks authored after the model's data cutoff; we have run none of them. Our
\codebleu implementation uses a node-type multiset overlap and an
identifier-level dataflow heuristic.

\textbf{Internal.} Design A is an inference-contract ablation on a frozen
checkpoint whose training corpus had no signature field, so the comparator prompt
was itself off-distribution; the effect size is far larger than that perturbation
plausibly explains, but the exact figures are conditional on that checkpoint. The
component isolations in Table~\ref{tab:sweep} are single-seed diagnostics. The
studies summarised in Section~\ref{sec:prior} are unreplicated.

\textbf{External.} One source language, two ISAs, one decoder family. The
frontier measurement rests on one provider's model on one date. Our survey
(Section~\ref{sec:scope}) revises what we previously believed about transfer:
HumanEval-Decompile is clean, but the same system's real-world splits ship the
true symbol in the assembly, so the critique \emph{does} reach C evaluation
practice --- through a channel we had not considered until the survey found it.
What remains open is whether the \emph{channel decomposition} reproduces on
another language and instruction set.
Section~\ref{sec:arm64} sets out exactly what that costs on our ARM64 corpus ---
inference alone on a frozen checkpoint, with the views already built and the
channels physically separated --- along with the one cell it cannot test. We
regard it as the most valuable next study rather than a limitation we can argue
away. We note that it would also vary the model architecture and the input
representation alongside the instruction set, which forecloses any comparison of
absolute magnitudes between the two corpora but makes a surviving dissociation a
conceptual replication rather than a retest.

\textbf{Conclusion.} Several nulls are underpowered by construction, and where
discordant counts are small we report the minimum attainable $p$ so a null is not
mistaken for evidence of absence. We apply Section~\ref{sec:turnover} by
discordance: aggregate comparisons and the $35$-to-$0$ identifier containment
stand on their exact $p$-values, while task-level claims resting on a handful of
discordant pairs --- the solved-set overlaps of Section~\ref{sec:nobinary} --- we
decline to make.

\section{Conclusion}
Two experiments, one an inference-time ablation and one an untrained baseline,
show that HumanEval-derived decompilation benchmarks
substantially measure their own semantic identifiers: removing the function name,
the parameter names and the disassembly symbol --- while holding types, arity and
the instruction stream fixed --- collapses functional correctness by a factor of
six while \emph{improving} compilation, and an untrained baseline given the
signature alone is not measurably worse than models fine-tuned on binaries.
Exploratory architecture, fine-tuning and objective-side sweeps move nothing
beyond the design's resolution. Against that, one channel effect --- the typed
contract on functional correctness --- is certified at $p{=}0.031$ under a
pre-registered six-seed test with a hard stop.

Those identifiers are a specification. Whether the model recalls a memorised
solution or reads the specification and solves the task afresh, our design cannot
say, and we have stopped claiming it can. Either way the benchmark cannot
attribute the score to the binary.

Underneath those results is the measurement problem of Section~\ref{sec:turnover}:
at these success rates a \passatk score replicates to one task while the tasks
behind it turn over almost completely, so the number certifies an aggregate, not
a capability on any particular function.

Any input channel teaches a small model the syntax and interface of a modern
language; none of the channels we tried teaches its semantics until the input
carries the information the compiler erased. The practical consequences are
narrow and firm: \emph{scrub the identifiers --- including the one in the
disassembly --- and report that arm}, or the number is not a measurement of
decompilation; and \emph{report the tasks solved in every seed beside the mean},
or the number is not a measurement of a capability. Across six published
evaluation configurations, no one currently does either.

\section*{Data Availability}
The replication package --- corpora and manifests, the object-pool extractor, the
sealed frontier harness with its hashes, per-seed generations and scores, the
pre-registration files, the contamination notice and the machine-readable
integrity reports --- is archived at
\mbox{\url{https://doi.org/10.5281/zenodo.21822791}} with code under Apache-2.0 and documentation under CC-BY-4.0; raw
benchmark-derived material remains under its upstream terms as recorded in the
package's licence map.

\bibliographystyle{IEEEtran}
\bibliography{thesis}

\end{document}
